\section{Independent signed slopes and smooth completion}
\label{sm:filling-proof}

Fix the subset $I$ and set all other $\chi_i$ equal to zero.
For $i\in I$ write $r_i=1/p_i$. Solve
\begin{equation}\label{sm:reciprocal}
                  V_i(\boldsymbol\chi)+r_iU_i(\boldsymbol\chi)
                                  =-2\pi\ii r_i.
\end{equation}
At $\mathbf r=0$ the solution is $\boldsymbol\chi=0$.
Equations \eqref{sm:UV} and \eqref{sm:Hessian} give the invertible
Jacobian
\begin{equation}\label{sm:Jacobian}
                       D_{\boldsymbol\chi}V(0)
                              =-2\pi\ii\diag(\tau_i)_{i\in I}.
\end{equation}
The implicit-function theorem gives a unique solution on one
reciprocal-slope polydisk. Shrink it until the representation and
positive original chain are in their common fixed neighborhood.
For real signed $p_i$ this imposes a lower bound on each $|p_i|$,
with no restriction on their ratios. For $p_i\ne0$ the equations are
\begin{equation}\label{sm:primitive-equations}
                           U_i+p_iV_i=-2\pi\ii .
\end{equation}
They are precisely the critical equations of
\eqref{sm:filled-action}.

The peripheral direction is actually loxodromic at these solutions.
Separate evenness makes the apparent quotient
$\widetilde\tau_i=V_i/U_i$ removable, holomorphic near zero, and
equal to $\tau_i$ there. Shrink once more so that
$\Ima\widetilde\tau_i>0$. Equation \eqref{sm:reciprocal} yields
\begin{equation}\label{sm:positive-length}
 \chi_{i,*}=\frac{r_i}{\widetilde\tau_i+r_i},\qquad
 \Rea V_i=\frac{2\pi r_i^2\Ima\widetilde\tau_i}
                         {|\widetilde\tau_i+r_i|^2}>0 .
\end{equation}
This is valid for either sign of a nonzero real reciprocal slope.
It is a classical consequence of the positive primitive cusp ratio.

We now apply the precise near-complete completion result
\cite[Proposition~6.1]{Marked}. Its input is an oriented complete
finite-volume manifold with marked primitive cusp bases, and an
actual nearby framed representation with small similarity logs in
bounded frames. At unfilled cusps both multipliers must be one.
At each filled cusp a primitive integral combination of the small
logs must equal $2\pi\ii$ or $-2\pi\ii$. It then completes the
original topological filling to a smooth complete hyperbolic metric.

Every hypothesis has been supplied on this enlarged syntax:
Theorem~\ref{sm:complete} constructs the actual complete exterior;
the unchanged two-rectangle collar proves its primitive basis;
Sections~\ref{sm:maps}--\ref{sm:group} construct the actual framed
representation and the original group relations; all unfilled
cusps have translation holonomy; and \eqref{sm:primitive-equations}
uses the primitive vector $(1,p_i)$. The fixed small parameter
neighborhood is chosen before the slopes.

\subsection{Why the integer equation gives one actual angular turn}
We recall the part of the imported proof that prevents an exponential
period equation from being mistaken for completion. On a complete
cusp link use a translation lattice $\Lambda\subset\CC$ and its fixed
generator paths. A nearby peripheral action is affine,
$A_\lambda(z)=a_\lambda z+b_\lambda$, in a bounded moving flag
frame. A smooth partition $\eta$ on the complete plane with
$\sum_{\lambda\in\Lambda}\eta(w-\lambda)=1$ gives
\[
 d_\rho(w)=\sum_{\lambda\in\Lambda}
                    \eta(w-\lambda)A_\lambda(w-\lambda).
\]
On a fundamental parallelogram this is $C^1$-close to the complete
identity and has positive real Jacobian; equivariance propagates
this fact. With the real peripheral logarithmic character in the
height coordinate it gives a nonsingular developed cusp, matched
over a compact collar to a nearby developing map on the original
truncated core.

At a filled cusp a multiplier has nonunit modulus by
\eqref{sm:positive-length}. Commutativity gives a second common
fixed point $q$ besides infinity. Since the original complete
translations are nonzero, $|q|\to\infty$ near completeness.
The developed fixed generator paths stay bounded, with bounded
length. They miss $q$, and their increments of
$f=\Log(d_\rho-q)$ tend to zero. Quantitatively, a path of length
at most $L$ contained in $|d_\rho|\le R$ has
$|\Delta f|\le L/(|q|-R)$.
Thus those increments are the prescribed small logs $U,V$,
without an added integral multiple of $2\pi\ii$, \emph{before}
multiplication by the large integer $p$.

The metric $|df|^2$ descends to a compact nonsingular flat torus.
Its complete lifted developing map is a diffeomorphism to $\CC$,
so $U,V$ form its actual rank-two period lattice. The developed
end contains a punctured metric tube with cylindrical metric
\begin{equation}\label{sm:tube}
             dR^2+\cosh^2R\,dX^2+\sinh^2R\,d\theta^2.
\end{equation}
A complex period $L$ acts by
$(X,\theta)\mapsto(X+\Rea L,\theta+\Ima L)$.
The original primitive slope has period $U+pV=-2\pi\ii$,
exactly one angular turn. Its complementary primitive period $V$
has nonzero real part. Adding $R=0$ produces the smooth solid
torus with this original meridian. The slope fixes both
$q,\infty$ and has multiplier one, so it is the identity in the
group, not a nontrivial unipotent with the same eigenvalues.
The remaining ends are full horoball quotients, and the middle
region is compact. Completeness and finite volume follow.

\begin{figure}[ht]
\centering\begin{tikzpicture}[x=.9cm,y=.9cm,>=Stealth,
 every node/.style={font=\small}]
 \coordinate (O) at (0,0); \coordinate (V) at (3,1.1);
 \coordinate (M) at (0,-2.3); \coordinate (U) at (-3,-3.4);
 \draw[gray!55,dashed] (V)--(3,-1.2)--(M);
 \fill[blue!7] (O)--(V)--(3,-1.2)--(M)--cycle;
 \draw[->,thick,blue!60!black] (O)--node[above] {$V$}(V);
 \draw[->,thick] (O)--node[right] {$m=-2\pi\ii$}(M);
 \draw[->,thick,orange!75!black] (O)--node[left] {$U=m-V$}(U);
 \draw[dashed,->,gray] (U)--node[below] {$V$}(M);
 \fill (O) circle (1.6pt); \node[above left] at (O) {$0$};
 \node[align=left,anchor=west,text width=4.2cm] at (4.8,-.7)
 {Original primitive basis: $U,V$.\\[3pt]
  Filled meridian: $m=U+pV$.\\[3pt]
  The fixed small increments are established before forming $m$.};
\end{tikzpicture}

\caption{The developed logarithmic link of a filled original
crossing circle. The old basis is $U,V$, while
$m=U+pV=-2\pi\ii$ is the primitive meridian of the added solid
torus. The picture shows $p=1$; the proof uses the small periods
before multiplying by any integer $p$. In the tube metric,
$m$ makes one angular turn and $V$ gives a nonzero axial
translation.}
\label{sm:period-figure}
\end{figure}

This verifies the metric conclusion of Theorem~\ref{sm:fillings}
for the actual disk-filled exterior. In particular high-index
estimates for a quantum coefficient, source positivity and whole
character-variety rigidity have played no role.

\section{Absolute volume on the original chain}\label{sm:absolute}

Set $G=2\pi\ii\Phi_\eps$. The exact peripheral identity gives
\begin{equation}\label{sm:primitive}
 \dd G=\frac12\sum_i V_i\,\dd U_i,\qquad
 \Ima G(0)=2\pi F_0=\Vol(Y_E^\eps).
\end{equation}
This is the $V\,dU$ polarization. A primitive of
$-\tfrac12\sum U\,dV$ differs by the boundary term
$\tfrac12\sum UV$; its regulator correction cannot be reused
unchanged here.

Let $\mathcal V(\rho)$ be the signed sum of the ideal tetrahedron
volumes of the retained original chain. The local boundary
differential \cite[equation~(7.3)]{Marked} reads, for similarity logs,
\begin{equation}\label{sm:volume-differential}
 \dd\mathcal V
     =\frac12\sum_i\bigl(\Rea V_i\,\dd\Ima U_i
                          -\Rea U_i\,\dd\Ima V_i\bigr).
\end{equation}
Its hypotheses hold on our positive nondegenerate chain:
internal faces cancel by the actual maps, edge-log sums stay fixed,
and the primitive cusp orientation is
$\Ima(\lambda/\mu)>0$ by \eqref{sm:tau}.
All main logs vanish. The proof of that local identity differentiates
Bloch--Wigner volume and cancels normal-face flows; it uses the
positive ordered model $(0,\infty,1,z)$ and the primitive cusp
intersection pairing. It is not a theorem restricted to the
complete packing of a particular diagram.

Taking the imaginary part of \eqref{sm:primitive} and differentiating
its boundary correction now gives
\begin{equation}\label{sm:regulator}
 \mathcal V(\rho_{\boldsymbol\chi})
      =2\pi\Rea\Phi_\eps(\boldsymbol\chi)
                     -\frac12\sum_i\Rea U_i\,\Ima V_i .
\end{equation}
The integration constant is fixed by the actual degree-one complete
cell volumes, not just by a derivative computation.
At a filling solution write $\chi_i=u_i+\ii v_i$. Since
$V_i=2\pi\ii(\chi_i-1)/p_i$,
\begin{equation}\label{sm:boundary-term}
 -\frac12\Rea U_i\,\Ima V_i
    =-\frac{2\pi^2}{p_i}v_i(u_i-1)
    =2\pi\Rea\left[\frac{\pi\ii}{2p_i}(\chi_i-1)^2\right].
\end{equation}
Unfilled auxiliaries have $U_i=V_i=0$.

\subsection{The same boundary triangulation after filling}
It remains to identify $\mathcal V$ with the volume of the smooth
metric just constructed. We verify the chain hypotheses of
\cite[Section~7.3]{Marked}. At a filled cusp conjugate the common
fixed points to $0,\infty$. On each triangle of the retained
boundary triangulation the ideal conical development is
$(a(w),Yb(w))$, with $b>0$. All functions and their derivatives
are bounded on a finite fundamental set. Choose an equivariant
axis map $(0,Y e^{\eta(w)})$, where $\eta$ has the real peripheral
logarithmic character. On the \emph{same} triangulation set
\begin{equation}\label{sm:chain-homotopy}
 H_s(w,Y)=\left(sa(w),
      Y\exp\bigl((1-s)\eta(w)+s\log b(w)\bigr)\right),
                  \qquad 0\le s\le1.
\end{equation}
Both horizontal coordinates and heights have the required
equivariance, and the homotopy agrees along every face.
Its height is comparable to $Y$, horizontal derivatives are bounded,
and height derivatives are $O(Y)$. The volume form
$y^{-3}\,dx_1\wedge dx_2\wedge dy$ therefore bounds its
three-dimensional boundary sweep by $O(Y^{-2})$.

Once the primitive slope is killed, extend that axis-valued boundary
over the added solid torus. The remaining cyclic group acts along
the axis, so all added vertices can lie on the axis and all added
tetrahedra have zero signed volume. At unfilled cusps the same
vanishing boundary estimate applies. Stokes on finite truncations,
followed by $Y\to\infty$, identifies the retained ideal chain with
the filled equivariant fundamental class. Comparing it with the
complete metric developing map on a common subdivision gives
\[
                 \mathcal V(\rho_{\boldsymbol\chi_*})
                        =\Vol(Y_E^\eps(\mathbf p;I)).
\]
An unrelated chain with the same peripheral differential would not
suffice; \eqref{sm:chain-homotopy} supplies the absolute comparison.
Combining it with \eqref{sm:regulator} and
\eqref{sm:boundary-term} proves \eqref{sm:filled-volume}.

\section{Proof boundary and finite algebra}\label{sm:scope}

The strict variational existence, Hessian and boundary-profile
identities are imported from \cite{Profiles}, with their precise
nonmetric content in Section~\ref{sm:variational}. The only local
volume input is the classical ideal-tetrahedron formula
\eqref{sm:tetra}. The local cells, all supporting planes and
cross ratios, entire rational prism maps, collar words and
similarity multipliers are proved here. The individual-mask
induction establishes their enlarged closed topology and every
actual versus completion-map status. The earlier
completed-root papers \cite{Geometry,Fillings} are archival
accounts of that narrower construction, not prerequisites of
the present local calculations or global theorem.
Smooth metric completion and the local
volume differential are the explicitly stated inputs from
\cite{Marked}, with all hypotheses verified here.

No source asymptotic, full spectral bound or nonzero Fourier
contribution is claimed. To identify the classical action with a
quantum packet one must independently match each endpoint-normalized
leaf action, including its color-dependent additive constant, and
then eliminate these same wires. A further whole-source estimate
and filling noncancellation theorem are needed for an invariant
volume limit. The geometry alone does not supply them.
Similarly no equivalence with a completed-root presentation is used,
and no moving main-meridian or preferred-longitude deformation
claim is made.

The accompanying \texttt{check.py} verifies the finite mask closure
and its allowed edge pairings, the literal endpoint-connection
table, projective turn/closing and pure/odd return identities,
the independent-phase prism table, complete collar conjugations,
deformed peripheral multipliers, and the primitive Hessian and
critical boundary normalization.
It includes altered mask, pairing and factor-of-two controls.
It writes a source-hashed JSON record with \texttt{--output}.
These exact finite checks do not prove convexity, the variational
existence theorem, the original manifold identification,
implicit-function theory, smooth completeness or the chain
comparison.

\begin{thebibliography}{9}
\bibitem{Profiles}
\emph{Stationary profiles of strict Newton trees},
companion paper, Lemma~2.2, Theorem~3.3, Lemma~4.1 and
Proposition~4.2.
\href{https://sashakolpakov.github.io/murakamis_dream/papers/strict-newton-profiles.pdf}
{Companion PDF}.

\bibitem{Geometry}
\emph{Pyramids, prisms, and complete geometry of Conway closures},
archival companion for the completed-root grammar.
\href{https://sashakolpakov.github.io/murakamis_dream/papers/conway-geometry.pdf}
{Companion PDF}.

\bibitem{Fillings}
\emph{Marked deformations and long fillings of Conway closures},
archival companion for the completed-root grammar.
\href{https://sashakolpakov.github.io/murakamis_dream/papers/conway-fillings.pdf}
{Companion PDF}.

\bibitem{Marked}
\emph{Circle patterns and primitive fillings of augmented links},
companion paper, Proposition~6.1, equation~(7.3) and Section~7.3.
\href{https://sashakolpakov.github.io/murakamis_dream/papers/marked-fillings.pdf}
{Companion PDF}.

\bibitem{Milnor}
J.~Milnor, Hyperbolic geometry: the first 150 years,
\emph{Bull.\ Amer.\ Math.\ Soc.\ (N.S.)} \textbf{6} (1982), 9--24.
\href{https://doi.org/10.1090/S0273-0979-1982-14958-8}
{doi:10.1090/S0273-0979-1982-14958-8}.

\bibitem{Purcell}
J.~S.~Purcell, An introduction to fully augmented links,
\emph{Contemp.\ Math.} \textbf{541} (2011), 205--220.
\href{https://users.monash.edu.au/~jpurcell/papers/expos-auglinks.pdf}
{Author's manuscript}.

\bibitem{NZ}
W.~D.~Neumann and D.~Zagier, Volumes of hyperbolic three-manifolds,
\emph{Topology} \textbf{24} (1985), 307--332.
\href{https://doi.org/10.1016/0040-9383(85)90004-7}
{doi:10.1016/0040-9383(85)90004-7}.
\end{thebibliography}
